\chapter[Protocolo experimental]{Protocolo experimental}
\label{chap: protocolo}

O \autoref{chap: metodologia} descreveu o que a busca faz. Este capítulo descreve como o seu resultado é medido e comparado: as réguas que traduzem a pergunta de pesquisa em números, os pontos de referência contra os quais essas réguas são lidas, os braços experimentais, a análise estatística e as validações auxiliares. Com exceção das varreduras de calibração, feitas antes da bateria (\autoref{sec: calibracao}), nada do que é descrito aqui altera a busca. As medições incidem sobre os elencos que ela devolve e sobre elencos de referência, sempre sob sorteios que a busca nunca viu.

\section{Das questões às réguas}
\label{sec: reguas}

A pergunta central do trabalho se desdobra nas três questões formuladas na \autoref{sec: questoes}. Cada uma é respondida por um conjunto de métricas, resumido no \autoref{quad: reguas}.

\begin{quadro}[htb]
    \caption{Questões de pesquisa, métricas que as respondem e relação de cada métrica com a aptidão.}
    \label{quad: reguas}
    \centering
    \small
    \begin{tabularx}{\textwidth}{l X l}
        \toprule
        \textbf{Questão} & \textbf{Métricas} & \textbf{Relação com a aptidão} \\
        \midrule
        QP1: equilíbrio & $\Pdom$ e seus três termos; número de \textit{hard counters}; personagens com taxa global em $[0{,}40;\, 0{,}60]$; elenco equilibrado (\autoref{eq: equilibrado}) & otimizadas \\
                        & Geração de convergência & registrada no laço \\
                        & Replicação e robustez a regras não vistas & \textit{post hoc} \\
        \midrule
        QP2: identidade & Estrutural: $\Pdrift$ & otimizada \\
                        & Estrutural: validador, Camadas 1 e 2 & parcialmente endógena \\
                        & Funcional: validador, Camada 3; concordância $\tau$ & \textit{post hoc}, \textit{held-out} \\
        \midrule
        QP3: compromisso & Fronteira de Pareto, hipervolume e espaçamento & \textit{post hoc} \\
                         & Relação de Pareto por semente; comparação entre braços & \textit{post hoc} \\
        \bottomrule
    \end{tabularx}
    \source
\end{quadro}

A identidade é medida por duas réguas, uma de cada lado da linha que separa premissa e resposta (\autoref{sub: premissa}). A régua \textbf{estrutural} pergunta se os genes continuam reconhecíveis. É o eixo que $\Pdrift$ otimiza, e as Camadas 1 e 2 do validador (\autoref{sub: validador}) medem o mesmo eixo por outro caminho. Por isso são \textbf{parcialmente endógenas}: um bom resultado nelas reflete, em parte, a penalidade de desvio ter funcionado. A régua \textbf{funcional} pergunta se cada personagem continua \textit{jogando} como o seu arquétipo. É ela que responde à pergunta de pesquisa, que fala em identidades funcionais. Nenhum termo da aptidão se refere a comportamento, e as métricas funcionais são, nesse sentido, \textit{held-out}.

Elas não são, contudo, independentes dos genes. Cada asserção comportamental é consequência quase direta de um gene definidor: o atordoamento infligido decorre do gene de atordoamento, o dano arrancado através da guarda decorre do agarrão, a distância média de luta decorre do alcance e do recuo. Um elenco pode preservar esses genes e perder o comportamento, porque o comportamento depende também do oponente, mas o contrário é pouco provável. A régua funcional é uma medida que a busca não vê, e não uma medida sem relação com o que a busca otimiza.

O ciclo de vantagens do \autoref{quad: ciclo_canonico} \textbf{não} é uma régua de identidade. Ele é uma hipótese autoral sobre quem deveria vencer quem, e o próprio elenco canônico a realiza apenas parcialmente no modelo (\autoref{chap: resultados}). O ciclo é tratado como hipótese a ser testada, uma única vez e na resolução que o teste exige (\autoref{sec: protocolo_ciclo}).

\section{Instrumentos de medição de identidade}
\label{sec: instrumentos}

\subsection{Perfil comportamental}
\label{sub: perfil}

O perfil comportamental de um personagem é a média, sobre os seus quatro confrontos, de dez métricas extraídas do registro das lutas sub-\textit{tick} a sub-\textit{tick} (\autoref{tab: perfil}). Ele é medido com 200 lutas por par, sob uma semente de avaliação fixa, a mesma para todos os elencos comparados. As frações de postura são calculadas sobre os sub-\textit{ticks} em que o personagem age, e a defesa é separada em escolhida e forçada, para que o encurralamento não se confunda com uma disposição defensiva (\autoref{sub: postura}).

\begin{table}[htb]
    \caption{Métricas do perfil comportamental de um personagem.}
    \label{tab: perfil}
    \centering
    \small
    \begin{tabularx}{\textwidth}{l X}
        \toprule
        \textbf{Métrica} & \textbf{Definição} \\
        \midrule
        Avanço, recuo & Fração dos sub-\textit{ticks} ativos em postura de avanço e de recuo \\
        Defesa escolhida & Fração dos sub-\textit{ticks} ativos em defesa por intenção GUARDA \\
        Defesa forçada & Fração dos sub-\textit{ticks} ativos em defesa por falta de espaço para recuar \\
        Fora de alcance & Fração dos sub-\textit{ticks} da luta em que o personagem age sem alcançar o oponente \\
        Atordoado & Fração dos sub-\textit{ticks} da luta em que o personagem está atordoado \\
        Golpes conectados & Número de golpes que atingem o oponente, por luta \\
        Distância média & Distância média entre os lutadores ao longo da luta \\
        Atordoamento infligido & Sub-\textit{ticks} de atordoamento aplicados ao oponente, por luta \\
        Quebra de guarda & Dano adicional arrancado pelo agarrão através da guarda do oponente, por luta \\
        \bottomrule
    \end{tabularx}
    \source
\end{table}

\subsection{Validador de identidade}
\label{sub: validador}

O validador verifica 23 asserções de ordenação, listadas no \autoref{quad: validador} e organizadas em três camadas. A \textbf{Camada 1} (13 asserções) compara personagens entre si nos genes: o \textit{Rushdown} deve ser o mais rápido do elenco, o \textit{Turtle} o de mais pontos de vida, e assim por diante. As asserções da Camada 1 coincidem com os genes definidores (\autoref{sub: canonicos}). A \textbf{Camada 2} (5 asserções) compara genes de um mesmo personagem, normalizados pela amplitude dos limites: o \textit{Zoner}, por exemplo, deve ter alcance relativamente maior que a velocidade. A \textbf{Camada 3} (5 asserções) compara personagens no perfil comportamental: cada arquétipo deve liderar o elenco na métrica que expressa o seu estilo.

\begin{quadro}[htbp]
    \caption{As 23 asserções do validador de identidade.}
    \label{quad: validador}
    \centering
    \small
    \begin{tabularx}{\textwidth}{p{2.7cm} l X}
        \toprule
        \textbf{Camada} & \textbf{Arquétipo} & \textbf{Asserção} \\
        \midrule
        \multirow{5}{=}{1: genes, entre personagens (13)}
                        & \textit{Rushdown} & maior velocidade; menor intervalo entre golpes; maior probabilidade de FRENTE \\
                        & \textit{Zoner} & maior alcance; maior repulsão; maior probabilidade de RECUAR \\
                        & \textit{Combo Master} & maior atordoamento \\
                        & \textit{Grappler} & maior dano; maior poder de agarrão \\
                        & \textit{Turtle} & menor velocidade; maior intervalo entre golpes; mais pontos de vida; maior probabilidade de GUARDA \\
        \midrule
        \multirow{4}{=}{2: genes, no mesmo personagem (5)}
                        & \textit{Zoner} & alcance maior que velocidade \\
                        & \textit{Rushdown} & velocidade maior que alcance \\
                        & \textit{Grappler} & pontos de vida maiores que velocidade \\
                        & \textit{Combo Master} & atordoamento maior que repulsão; velocidade maior que alcance \\
        \midrule
        \multirow{5}{=}{3: comportamento (5)}
                        & \textit{Zoner} & maior distância média de luta \\
                        & \textit{Rushdown} & mais golpes conectados por luta \\
                        & \textit{Turtle} & maior fração de defesa escolhida \\
                        & \textit{Combo Master} & maior atordoamento infligido \\
                        & \textit{Grappler} & maior dano arrancado através da guarda \\
        \bottomrule
    \end{tabularx}
    \source
\end{quadro}

Duas regras impedem que o validador conceda asserções sem identidade. A primeira conta os \textbf{empates contra a asserção}: um personagem empatado com outro só é considerado ``o maior'' se o for estritamente. Sem essa regra, cinco cópias idênticas do mesmo personagem passariam em 4 das 13 asserções da Camada 1, apenas pela ordem em que aparecem no elenco. A segunda compara os pesos comportamentais pelas \textbf{probabilidades de intenção} da \autoref{eq: intencao}, e não pelos valores crus, pela mesma razão que motiva a reescala da \autoref{eq: reescala}: a escala dos pesos não tem efeito no combate e não pode alterar o veredito.

O validador é uma régua binária e grosseira. Cada asserção é um bit, e ficar em segundo lugar por pouco conta o mesmo que ficar em último. A Camada 3, em particular, tem apenas cinco bits. Por isso a identidade funcional é medida também por uma régua contínua.

\subsection{Concordância de ranking comportamental}
\label{sub: tau}

A concordância de ranking pergunta se a \textit{ordem} dos cinco personagens em cada métrica comportamental continua sendo a do elenco canônico, e não apenas quem lidera cada uma. Para cada métrica $m$ da \autoref{tab: perfil}, calcula-se o $\tau_b$ de Kendall \cite{kendall1938new, kendall1945treatment}, a versão do coeficiente corrigida para empates, entre o vetor dos cinco valores no elenco canônico e o vetor dos cinco valores no elenco medido. A concordância é a média sobre as dez métricas:
\begin{equation}
    \label{eq: tau}
    \tau(\mathbf{x}) = \frac{1}{10} \sum_{m=1}^{10} \tau_b\!\left( \mathbf{v}_m(\mathbf{c}),\, \mathbf{v}_m(\mathbf{x}) \right),
\end{equation}
em que $\mathbf{v}_m(\mathbf{x}) \in \mathbb{R}^5$ reúne os valores da métrica $m$ para os cinco personagens do elenco $\mathbf{x}$. O valor 1 indica a ordem canônica em todas as métricas; 0 indica ausência de relação; $-1$, a ordem invertida. O perfil canônico de referência é medido sob as mesmas condições do elenco comparado, isto é, o mesmo número de lutas e a mesma semente. Como a medida é relativa ao elenco nos dois lados, um deslocamento que afete todos os personagens igualmente, como lutas mais longas, não conta como perda de identidade.

A concordância é uma estimativa ruidosa, e o número de lutas foi escolhido por reteste. Medida três vezes sobre o mesmo elenco evoluído, sob sementes diferentes, ela variou de \tauRetesteCentoVinte{} com 120 lutas por par, e de \tauRetesteDuzentas{} com 200. A medição usa 200.

\subsection{Diferenciação}
\label{sub: diferenciacao}

A diferenciação mede quão distintos os cinco personagens de um elenco são entre si. Cada personagem é representado por um vetor com os oito atributos normalizados pela amplitude dos limites e as três probabilidades de intenção. A diferenciação do elenco é a distância euclidiana média entre os $\binom{5}{2} = 10$ pares de personagens. Ela vale zero num elenco de cinco cópias idênticas. É uma leitura descritiva, fora da família de testes: um elenco sorteado também é bem diferenciado sem ter identidade alguma, e a métrica não distingue as duas coisas.

\section{Modelos nulos}
\label{sec: nulos}

Nenhuma das métricas de identidade tem piso igual a zero. Um elenco sorteado ao acaso satisfaz, em média, cerca de \nuloValMedio{} das 23 asserções do validador, porque cada asserção de ordenação tem chance não nula de ser verdadeira por acaso. Uma leitura como ``o elenco evoluído satisfaz 17 asserções'' só tem significado diante desse piso. O trabalho mede, por isso, cada métrica sobre 35 \textbf{modelos nulos}, elencos que não passaram por nenhum processo de otimização:

\begin{itemize}
    \item cinco \textbf{espelhos}, cada um formado por cinco cópias do perfil canônico de um arquétipo. O espelho é a solução trivial do equilíbrio: perfeitamente simétrico e sem identidade alguma. É o ponto de referência para ``tão equilibrado quanto a simetria perfeita'';
    \item trinta elencos \textbf{aleatórios}, sorteados uniformemente dentro dos limites. São o piso absoluto: nem equilibrados, nem com identidade.
\end{itemize}

Cada métrica de um elenco é reportada com duas leituras. A \textbf{posição} localiza o valor entre um piso e um teto:
\begin{equation}
    \label{eq: posicao}
    \text{posição} = \frac{\text{valor} - \text{piso}}{\text{teto} - \text{piso}},
\end{equation}
em que o piso é a média dos 35 nulos e o teto é o elenco canônico, para as métricas de identidade, ou a média dos espelhos, para $\Pdom$. O \textbf{$p$-valor empírico} é a fração dos 35 nulos que igualam ou superam o valor observado. Com 35 nulos, sua resolução é $1/35 \approx 0{,}03$.

Os modelos nulos não controlam o \textit{método}. Eles não são otimizados, e superá-los em $\Pdrift$ é garantido por construção a qualquer busca que tenha esse termo na aptidão. O controle que isola o efeito do termo de identidade é outro, descrito na \autoref{sec: desenho}. Os nulos servem para dar escala às métricas: dizer o que é acaso em cada uma.

\section{Desenho experimental}
\label{sec: desenho}

\subsection{Braços}
\label{sub: bracos}

A bateria experimental compara cinco braços, listados na \autoref{tab: bracos}. Os dois primeiros são os algoritmos que a pergunta de pesquisa compara. Os dois seguintes são \textbf{controles}: repetem o algoritmo escalar com um único fator de desenho alterado, com a mesma amostra e o mesmo orçamento. O quinto é o híbrido, reportado como contribuição de método.

\begin{table}[htb]
    \caption[Braços da bateria experimental]{Braços da bateria experimental. Todos usam 20 execuções independentes (sementes 42 a 61), população de 300 e 150 gerações.}
    \label{tab: bracos}
    \centering
    \small
    \begin{tabularx}{\textwidth}{l X}
        \toprule
        \textbf{Braço} & \textbf{Papel} \\
        \midrule
        AG escalar & Algoritmo principal, na configuração da \autoref{tab: parametros} \\
        NSGA-II & Algoritmo de comparação; entra nos testes pelo representante \textit{scalar\_optimum} e, secundariamente, pelo \textit{best\_dominance} \\
        Controle $\lambda_{\mathrm{drift}} = 0$ & AG escalar sem o termo de identidade: o contrafactual da pergunta, equilibrar sem pagar pela identidade \\
        Controle sem semente & AG escalar com população inicial inteiramente aleatória, como a do NSGA-II: separa algoritmo de inicialização \\
        Híbrido & 75 gerações de NSGA-II seguidas de 75 de AG escalar, entregando a fronteira inteira (\autoref{sec: hibrido}) \\
        \bottomrule
    \end{tabularx}
    \source
\end{table}

O controle $\lambda_{\mathrm{drift}} = 0$ é o experimento que responde à QP2. Os modelos nulos mostram o que é acaso numa métrica; o controle mostra o que a busca produz \textit{sem} o termo de identidade, na mesma amostra e no mesmo orçamento. A diferença entre o algoritmo escalar e esse controle é o efeito do termo de identidade, isolado de todo o resto do método.

\subsection{Amostra e pareamento}
\label{sub: amostra}

Um algoritmo evolutivo é estocástico, e uma execução é uma amostra, não um resultado \cite{eiben2015introduction}. Cada braço é executado com 20 sementes. O tamanho foi escolhido por uma simulação de poder, com duas distribuições separadas por um efeito grande ($\hat{A}_{12} = 0{,}80$): com 20 execuções por braço, o efeito é detectado com probabilidade de \poderNVinte{}, contra \poderNDez{} com 10. A simulação foi feita quando o teste era não pareado e a correção cobria três métricas. O teste pareado adotado depois tende a ter mais poder nesse desenho, e a família atual, de até sete métricas, impõe uma correção mais severa. As duas mudanças atuam em sentidos opostos, e o número deve ser lido como ordem de grandeza, não como poder exato.

Os braços usam as \textbf{mesmas} 20 sementes, e a semente controla tudo o que é aleatório numa execução: a população inicial, os sorteios dos operadores e o \textit{stream} de avaliação de cada geração. A execução de semente $r$ de um braço e a de semente $r$ de outro formam, portanto, um bloco experimental, e não duas amostras independentes. O desenho é \textbf{pareado}, e a análise estatística o respeita.

\subsection{Reavaliação independente}
\label{sub: reavaliacao}

O elenco que uma execução devolve foi selecionado \textit{porque} teve boa avaliação sob os sorteios da última geração. Parte dessa avaliação é sorte, e relatá-la seria relatar a sorte selecionada. Todo elenco devolvido é, por isso, reavaliado fora do laço, sob uma semente de avaliação comum a todas as execuções de todos os braços, com 1000 lutas por par. A semente comum faz com que diferenças entre execuções reflitam os elencos, e não os sorteios da medição. As métricas de identidade funcional são medidas sob a mesma semente, com 200 lutas por par (\autoref{sub: perfil}).

O número de lutas da reavaliação foi escolhido pelo ruído de medição. A 1000 lutas por par, o desvio-padrão do $\Pdom$ de um mesmo elenco entre sorteios cai para \ruidoDomDpMil{}, contra \ruidoDomDpTreino{} a 150 (\autoref{sec: ruido}). Há também um viés. Como $\Pdom$ é função de $|p - 1/2|$, o ruído de medição o \textit{aumenta} em média: os mesmos elencos medidos a 150 lutas têm $\Pdom$ médio de \ruidoDomMediaTreino{}, e a 1000, de \ruidoDomMediaMil{}. Uma régua grosseira, portanto, subestima o equilíbrio, e tanto mais quanto mais equilibrado é o elenco, porque o viés é máximo quando as taxas de vitória verdadeiras estão perto de 50\%.

\subsection{Famílias de sementes}
\label{sub: sementes}

Cada uso de números aleatórios no protocolo tem a sua família de sementes, e as famílias são disjuntas. Nenhuma avaliação reutiliza os sorteios de outra, e a amostra que escolhe uma configuração não é a mesma que a avalia. O \autoref{quad: sementes} lista as famílias.

\begin{quadro}[htb]
    \caption{Famílias de sementes do protocolo experimental.}
    \label{quad: sementes}
    \centering
    \small
    \begin{tabularx}{\textwidth}{l X}
        \toprule
        \textbf{Sementes} & \textbf{Uso} \\
        \midrule
        42 a 61 & Execuções da bateria; geram os \textit{streams} de treino a partir de $42\,000$ \\
        1000 a 1004 & Execuções das varreduras de calibração; \textit{streams} de treino a partir de $1\,000\,000$ \\
        \textit{stream} da geração $+\ 100\,000$ & Confirmação da convergência, inédita a cada geração \\
        9999 & Reavaliação independente, medição de identidade e sorteio dos 30 elencos aleatórios \\
        $10\,000$ a $10\,009$ & Validação externa \\
        $950\,000 + 53\,k$, $k = 0, \ldots, 15$ & Teste do ciclo autoral (16 \textit{streams}) \\
        $42 + 100\,i + j$ e deslocamentos de $10\,000$ & Análise de sensibilidade ($i$: personagem; $j$: gene) e o seu piso de ruído \\
        $500\,000$ em diante & Medições avulsas do ruído de $\Pdom$, feitas na preparação do texto \\
        \bottomrule
    \end{tabularx}
    \source
\end{quadro}

\section{Análise estatística}
\label{sec: estatistica}

Cada comparação opõe o algoritmo escalar a um dos outros quatro braços. Seguindo as recomendações para a comparação de algoritmos estocásticos \cite{derrac2011practical, arcuri2011practical}, cada comparação combina um teste de hipótese, um tamanho de efeito e uma correção para múltiplas comparações.

\textbf{Teste.} O teste principal é o de postos sinalizados de Wilcoxon \cite{wilcoxon1945individual}, pareado e bicaudal, aplicado às 20 diferenças entre execuções de mesma semente. Pares com diferença nula são descartados, pois não informam direção. O teste é não paramétrico e não supõe normalidade, o que o torna adequado a amostras pequenas e a métricas limitadas, algumas delas inteiras. O teste U de Mann-Whitney \cite{mann1947test}, não pareado, é reportado ao lado, sem correção, para que se possa verificar que nenhuma conclusão depende da escolha do teste. Ignorar o pareamento seria conservador, mas não inválido.

\textbf{Tamanho de efeito.} Um $p$-valor pequeno indica que uma diferença é improvável sob a hipótese nula, não que ela seja grande. O tamanho de efeito é o $\hat{A}_{12}$ de \citeonline{vargha2000critique}, a probabilidade de uma execução do primeiro braço produzir valor maior que uma do segundo, contando empates pela metade:
\begin{equation}
    \label{eq: a12}
    \hat{A}_{12} = \frac{1}{n^2} \sum_{i=1}^{n} \sum_{j=1}^{n} \left( \mathbb{1}[a_i > b_j] + \tfrac{1}{2}\, \mathbb{1}[a_i = b_j] \right).
\end{equation}
O valor 0,5 indica ausência de efeito. As magnitudes seguem os limiares dos autores sobre $|\hat{A}_{12} - 0{,}5|$: desprezível abaixo de 0,06, pequeno até 0,14, médio até 0,21 e grande acima disso.

\textbf{Família e correção.} Testar várias métricas ao nível $\alpha = 0{,}05$ inflaria a chance de ao menos um falso positivo. Os $p$-valores de cada comparação são corrigidos pelo procedimento de \citeonline{holm1979simple} sobre uma \textbf{família fixa} de sete métricas, a mesma em todas as comparações: três de equilíbrio ($\Pdom$, número de \textit{hard counters} e número de personagens com taxa global em banda) e quatro de identidade ($\Pdrift$, Camadas 1 e 2 do validador, Camada 3 e concordância $\tau$). Fixar a família antes da bateria impede que as métricas sejam escolhidas depois de vistos os resultados. Uma única regra, também fixada antes, retira uma métrica da família: se as 40 observações da comparação forem todas iguais, não há teste definido, e a métrica é reportada apenas descritivamente.

\textbf{Leituras descritivas.} Duas leituras ficam fora da família de testes, para não encarecer a correção. A primeira é a decomposição de $\Pdom$ nos seus três termos, que mostra de onde vem uma diferença de equilíbrio (\autoref{sub: dominancia}). A segunda é a \textbf{relação de Pareto por semente}: o ponto $(\Pdom, \Pdrift)$ do algoritmo escalar é comparado com a fronteira inteira do NSGA-II da mesma semente, ambos medidos sob o mesmo \textit{stream}, o da última geração. Em cada semente, ou algum ponto da fronteira domina o do escalar, ou o do escalar domina algum ponto da fronteira, ou nenhum dos dois. A relação não depende da escolha de um representante.

\section{Qualidade da fronteira}
\label{sec: qualidade_fronteira}

Cada fronteira do NSGA-II é resumida por dois indicadores \cite{deb2001multiobjective}. O \textbf{hipervolume} \cite{zitzler1999multiobjective} é a área do espaço de objetivos dominada pela fronteira e limitada por um ponto de referência. Mede ao mesmo tempo a proximidade do ótimo e a extensão da fronteira, e é maior quanto melhor. O ponto de referência $(1{,}3;\, 0{,}4)$ foi ancorado nos modelos nulos: $\Pdom$ próximo ao do elenco canônico, que é o equilíbrio de partida, e $\Pdrift$ próximo ao dos espelhos, que é a identidade nula. Uma solução pior que isso em algum dos objetivos não contribui para o hipervolume, que é o que ela vale. A referência é fixa, para que o indicador seja comparável entre execuções. Com a referência nos máximos teóricos, $(2;\, 1)$, o hipervolume ocupava 90\% da área disponível em qualquer fronteira e mal distinguia uma da outra. O \textbf{espaçamento} \cite{schott1995fault} é o desvio-padrão das distâncias de cada ponto ao vizinho mais próximo e mede a uniformidade da distribuição dos pontos: é menor quanto mais uniforme.

\section{Estrutura de vantagens}
\label{sec: protocolo_ciclo}

O ciclo de vantagens do \autoref{quad: ciclo_canonico} é testado como hipótese, numa medição única e de alta resolução. A resolução é o ponto central. Num elenco equilibrado as taxas de vitória dos pares ficam próximas de 50\%, e a margem mediana de uma aresta nos elencos do algoritmo escalar é de apenas 0,048. A 200 lutas por par, o desvio-padrão binomial é 0,035, e a direção de uma aresta seria, na prática, sorteada. O teste usa, por isso, 16 \textit{streams} de 1000 lutas por par, ou 16\,000 lutas por aresta, com desvio-padrão de 0,004.

Uma aresta é considerada \textbf{decidida} quando a taxa de vitória do par se afasta de 50\% por mais de dois desvios-padrão, e \textbf{mantida} quando é decidida na direção autoral. Arestas indecisas são excluídas da contagem, porque contá-las misturaria ruído nos dois sentidos. A hipótese nula é que o elenco não tem relação com o ciclo, de modo que cada aresta decidida segue a direção autoral com probabilidade 1/2. Ela é testada de duas formas: por um teste binomial sobre o total de arestas mantidas entre as decididas, e pela comparação das taxas de manutenção dos elencos de cada braço com as dos 30 elencos aleatórios, por Mann-Whitney e $\hat{A}_{12}$. A comparação usa a taxa, e não a contagem, porque um elenco aleatório decide quase todas as arestas e um equilibrado decide menos, e a contagem crua puniria o equilibrado por ter uma aresta sobre o limiar.

Ao lado do ciclo autoral, é medida a estrutura de vantagens sem referência a ele: o número de \textbf{tríades circulares} do torneio \cite{kendall1940method}, isto é, de trios em que A vence B, B vence C e C vence A. O valor é 0 numa ordem estrita, 2,5 em média num torneio em que cada confronto tem vencedor sorteado e 5 no máximo, quando cada personagem vence dois e perde dois. Um elenco sorteado dentro dos limites não se comporta como esse torneio: um personagem forte tende a vencer todos, e esses elencos costumam ser quase hierárquicos. A medida deve ser lida com cautela, porque um elenco globalmente equilibrado com pares decididos é forçado a ser intransitivo: uma ordem estrita entre cinco personagens teria taxas globais de 100\%, 75\%, 50\%, 25\% e 0\%.

\section{Validação externa}
\label{sec: validacao_externa}

A validação externa, no espírito da avaliação independente da aptidão empregada por \citeonline{browne2010evolutionary}, faz duas perguntas distintas a um elenco já evoluído.

A \textbf{replicação} pergunta se o equilíbrio medido durante a busca se confirma com uma amostra grande e sorteios inéditos, sob as mesmas regras de combate do treino. A \textbf{robustez} pergunta se o equilíbrio sobrevive a regras que a busca nunca viu. Trocar apenas a semente repetiria a primeira pergunta com mais amostra; mudar a regra testa se o equilíbrio depende das condições exatas em que foi otimizado. São perturbadas, uma de cada vez, quatro regras de combate: a distância inicial (40 e 60, em vez de 50), o comprimento do campo (80 e 120, em vez de 100), a persistência da intenção (4 e 6, em vez de 5) e o multiplicador da guarda (0,55 e 0,65, em vez de 0,6). Cada perturbação permanece dentro do que o modelo sustenta: a distância inicial continua maior que qualquer alcance, e o multiplicador da guarda desloca o ponto neutro do agarrão em apenas 0,05.

Em cada condição, 10 sementes de 500 lutas por par formam uma amostra de 5000 lutas por par. Cada taxa de vitória recebe um intervalo de confiança de Wilson a 95\% \cite{wilson1927probable} e é classificada contra a sua faixa de equilíbrio, $[0{,}40;\, 0{,}60]$ para a taxa global de cada personagem e $[0{,}35;\, 0{,}65]$ para cada par: \textbf{dentro}, se o intervalo inteiro está dentro da faixa; \textbf{fora}, se está inteiro fora; \textbf{inconclusiva}, se cruza a borda. A condição é \textbf{robusta} se todas as 15 taxas estão dentro, \textbf{frágil} se alguma está fora, e \textbf{inconclusiva} nos demais casos. O veredito não se torna mais severo com o número de sementes, como aconteceria com uma regra do tipo ``falhou em alguma das $K$ sementes''.

A validação externa é aplicada ao elenco canônico e a três elencos da execução de semente 42: o do algoritmo escalar e os representantes \textit{scalar\_optimum} e \textit{knee\_point} do NSGA-II.

\section{Análise de sensibilidade}
\label{sec: sensibilidade}

A análise de sensibilidade pergunta se a busca \textit{enxerga} cada gene, isto é, se um passo típico de mutação nesse gene move o desfecho das lutas. Ela é feita sobre um elenco equilibrado, o do algoritmo escalar de semente 42. Medi-la no elenco canônico seria inútil: com taxas de vitória presas em 0\% e 100\%, nenhuma perturbação muda o desfecho, e todo gene pareceria neutro por efeito de teto.

Para cada gene de cada personagem, o gene é deslocado numa janela de largura $2\sigma$, em que $\sigma$ é o desvio-padrão da mutação daquele gene (\autoref{sec: ag_escalar}), e mede-se a variação da taxa de vitória global do personagem entre os dois extremos da janela. Se a janela ultrapassa um limite do gene, ela é deslizada para dentro do intervalo, mantendo a largura, para que genes próximos a um limite não pareçam menos sensíveis apenas por isso. Os dois extremos são avaliados sob os mesmos sorteios, com 200 lutas por par. A estatística de cada gene é a média, sobre os cinco personagens, da variação absoluta.

O piso de ruído é medido na mesma estatística, sob a hipótese nula: uma janela de largura zero avaliada sob sorteios diferentes nos dois extremos, em três repetições, das quais se toma o maior valor. Um gene é \textbf{neutro} se a sua estatística não excede o piso, \textbf{limítrofe} se não excede o dobro do piso, e \textbf{visível} acima disso.

\section{Calibração}
\label{sec: calibracao}

Quatro dos parâmetros marcados como medidos na \autoref{tab: parametros} foram escolhidos por varreduras realizadas antes da bateria: a razão $\lambda_{\mathrm{drift}}/\lambda_{\mathrm{dom}}$, os pesos dos termos de $\Pdom$, a taxa de elitismo e o tamanho do torneio. Os demais foram medidos nas condições descritas junto de cada um, no \autoref{chap: metodologia}. As varreduras usaram orçamento reduzido: população de 120, 60 gerações e 5 sementes (1000 a 1004), disjuntas das da bateria. Cada braço de varredura altera um único parâmetro do algoritmo escalar, e os elencos resultantes são reavaliados como na bateria, com 1000 lutas por par. A \autoref{tab: calibracao} resume os braços. Ela traz o termo global $\tglob$ em todos eles e $\Pdom$ apenas onde os seus pesos não variam, pois nos braços que os variam o valor de $\Pdom$ não é comparável.

\begin{table}[htbp]
    \caption[Varreduras de calibração do algoritmo escalar]{Varreduras de calibração do algoritmo escalar (5 sementes, população 120, 60 gerações). Médias sobre as sementes; na primeira linha, a configuração adotada.}
    \label{tab: calibracao}
    \centering
    \footnotesize
    % Fonte: results/exploratory/multi_run_ga_pop120_gen60_n5_seed1000*.json
    \begin{tabularx}{\textwidth}{X c c c c c c}
        \toprule
        \textbf{Braço} & $\Pdom$ & $\tglob$ & \textbf{\textit{Hard counters}} & \textbf{Equilibrados} & $\Pdrift$ & $\tau$ \\
        \midrule
        \textbf{Adotada}: $\lambda_{\mathrm{drift}} = 1$; $\beta = 1/0{,}5/0{,}5$; elitismo 10\%; torneio 3 & 0,042 & 0,040 & 0,4 & 3/5 & 0,245 & 0,42 \\
        \midrule
        $\lambda_{\mathrm{drift}} = 0{,}25$ & 0,041 & 0,032 & 1,2 & 2/5 & 0,377 & 0,02 \\
        $\lambda_{\mathrm{drift}} = 0{,}5$  & 0,044 & 0,039 & 0,4 & 3/5 & 0,346 & $-$0,05 \\
        $\lambda_{\mathrm{drift}} = 2$      & 0,183 & 0,042 & 4,2 & 1/5 & 0,139 & 0,60 \\
        $\lambda_{\mathrm{drift}} = 4$      & 0,356 & 0,069 & 7,8 & 0/5 & 0,090 & 0,72 \\
        \midrule
        $\beta = 1/0/0$       & -- & 0,023 & 8,6 & 0/5 & 0,156 & 0,53 \\
        $\beta = 1/0{,}5/0$   & -- & 0,031 & 2,6 & 1/5 & 0,248 & 0,30 \\
        $\beta = 1/1/1$       & -- & 0,030 & 0,0 & 5/5 & 0,301 & 0,11 \\
        $\beta = 1/2/0{,}5$   & -- & 0,041 & 0,0 & 5/5 & 0,348 & 0,16 \\
        \midrule
        Elitismo 0\%   & 0,068 & 0,064 & 0,8 & 3/5 & 0,269 & 0,30 \\
        Elitismo 5\%   & 0,042 & 0,038 & 0,6 & 2/5 & 0,289 & 0,24 \\
        Elitismo 20\%  & 0,044 & 0,040 & 0,6 & 3/5 & 0,281 & 0,19 \\
        Elitismo 30\%  & 0,073 & 0,032 & 1,6 & 2/5 & 0,285 & 0,20 \\
        \midrule
        Torneio 2 & 0,040 & 0,040 & 0,2 & 4/5 & 0,264 & 0,36 \\
        Torneio 5 & 0,039 & 0,031 & 0,6 & 3/5 & 0,291 & 0,17 \\
        Torneio 7 & 0,038 & 0,032 & 0,6 & 3/5 & 0,340 & 0,05 \\
        \bottomrule
    \end{tabularx}
    \source
\end{table}

A varredura da razão $\lambda_{\mathrm{drift}}/\lambda_{\mathrm{dom}}$ mostra o joelho descrito na \autoref{sub: formulacoes}. Abaixo de 1, o equilíbrio não melhora e a identidade se perde: com $\lambda_{\mathrm{drift}} = 0{,}25$ e $0{,}5$, a concordância $\tau$ cai ao nível do acaso, e $\Pdom$ fica entre 0,041 e 0,044 nos três braços. Acima de 1, a identidade melhora, mas o equilíbrio se degrada: $\Pdom$ sobe para 0,183 e 0,356, e os \textit{hard counters}, de 0,4 para 4,2 e 7,8 por elenco. Entre os valores testados, a razão unitária é o maior peso de identidade que ainda não custa equilíbrio.

A varredura dos pesos de $\Pdom$ sustenta a presença dos termos secundários e expõe o compromisso da sua repartição, discutidos na \autoref{sub: dominancia}. As varreduras de elitismo e torneio mostram que nenhuma alternativa domina a configuração adotada. O torneio de tamanho 2 obtém menos \textit{hard counters} e mais elencos equilibrados, mas paga em $\Pdrift$ e em $\tau$; as demais alternativas pioram a identidade sem melhorar o equilíbrio de forma consistente. A leitura correta é ``testados, e nenhuma alternativa os supera em todas as réguas'', e não ``ótimos''.

Duas ressalvas limitam o uso dessas varreduras. Com 5 sementes, diferenças pequenas não se separam do ruído, e a tabela serve para ordenar configurações, não para declarar vencedores. Além disso, o orçamento reduzido não transfere necessariamente para o da bateria: a ordem entre o algoritmo escalar e o NSGA-II, por exemplo, se inverte com população de 120, razão pela qual nenhuma comparação entre algoritmos é feita em orçamento reduzido. A escolha do híbrido seguiu um critério escrito antes da varredura (\autoref{sec: hibrido}): eliminar as configurações consistentemente piores em equilíbrio e, entre as restantes, preferir a que melhor preserva identidade. Nenhuma das seis configurações sobreviveu à eliminação no orçamento reduzido, e a configuração levada à bateria foi a mais simples. Nesse orçamento, ela preservou menos identidade que o próprio algoritmo escalar (desvio mediano de \hibReduzidoDrift{} contra \agReduzidoDrift{}), de modo que o seu efeito só pôde ser avaliado na bateria.

\section{Reprodutibilidade e rastreabilidade}
\label{sec: reprodutibilidade}

Toda execução é reprodutível a partir da sua semente. A semente fixa a população inicial, os sorteios dos operadores e os \textit{streams} de avaliação, e cada luta é semeada pela sua posição na avaliação, de modo que o resultado não depende da ordem nem do número de processos que a executam.

Cada artefato gravado pelo sistema, dos resultados de uma execução às comparações estatísticas, abre com um registro de proveniência: o valor de cada parâmetro de configuração, um resumo criptográfico dos perfis canônicos, um resumo do código-fonte do motor e, nos artefatos de medição, um resumo do código de medição que os produziu. Ao carregar um artefato, o sistema compara esse registro com o estado atual e avisa o que mudou. Uma ferramenta que \textit{grava} um artefato novo a partir de outro recusa entrada desatualizada, pois o artefato novo seria carimbado com a configuração atual e passaria a atestar um número produzido por outra. Os números da bateria citados nos capítulos seguintes saem de artefatos gerados em 23 de setembro de 2026, sobre a configuração descrita no \autoref{chap: metodologia}. As poucas medições avulsas feitas na preparação do texto, como a duração das lutas e o ruído de medição de $\Pdom$, não têm artefato próprio e são identificadas como tais onde aparecem.
